%% Đề 15 -- de-kiem-tra-cuoi-hoc-ki-1-toan-9-nam-2022-2023-so-gddt-lang-son.pdf
\deheader{ĐỀ SỐ 15}%
  {Đề kiểm tra Học kỳ — Năm học 2022--2023}%
  {Sở GD\&ĐT Lạng Sơn | Chỉ phần Đại số}
\addcontentsline{toc}{section}{Đề 15 — Đề kiểm tra Học kỳ — Năm học 2022--2023}

\begin{de_cau}{Câu 1}{2,0 điểm}
\begin{enumerate}[label=\alph*)]
\item Tính: $\sqrt{20}-\sqrt{45}+3\sqrt{80}$.
\item Rút gọn: $M=\left(\dfrac{\sqrt{x}}{\sqrt{x}+3}+\dfrac{3}{\sqrt{x}-3}\right)\cdot\dfrac{\sqrt{x}-3}{2}$, $x\geq 0, x\neq 9$.
\end{enumerate}
\end{de_cau}

\begin{loi_giai_cau}
\textbf{a)} $2\sqrt{5}-3\sqrt{5}+12\sqrt{5}=11\sqrt{5}$.

\textbf{b)}
\[M=\frac{\sqrt{x}(\sqrt{x}-3)+3(\sqrt{x}+3)}{(\sqrt{x}+3)(\sqrt{x}-3)}\cdot\frac{\sqrt{x}-3}{2}
=\frac{x+6}{x-9}\cdot\frac{\sqrt{x}-3}{2}=\frac{(x+6)}{2(\sqrt{x}+3)}.\]
\end{loi_giai_cau}

\begin{de_cau}{Câu 2}{2,5 điểm}
\begin{enumerate}[label=\alph*)]
\item Giải phương trình: $3x^2+2x-1=0$.
\item Giải hệ phương trình: 
\[ \begin{cases}\dfrac{x}{2}+\dfrac{y}{3}=1\\\dfrac{x}{3}-\dfrac{y}{2}=\dfrac{5}{6}\end{cases} \].
\item Cho parabol $(P)\colon y=2x^2$ và đường thẳng $(d)\colon y=3x+2$. Tìm tọa độ giao điểm.
\end{enumerate}
\end{de_cau}

\begin{loi_giai_cau}
\textbf{a)} $\Delta=4+12=16$. $x=\dfrac{-2\pm 4}{6}$: $x_1=\dfrac{1}{3}, x_2=-1$.

\textbf{b)} Đặt $u=\dfrac{x}{6}, v=\dfrac{y}{6}$... Nhân (1) $\times 6$: $3x+2y=6$; (2) $\times 6$: $2x-3y=5$.

$(2)\times 3+(1)\times 2$: $4x+4x+\ldots$ Cộng nhân chéo: $9x+4x=18+10\Rightarrow x=2, y=0$.

\textbf{c)} $2x^2=3x+2\Rightarrow 2x^2-3x-2=0\Rightarrow(2x+1)(x-2)=0$.
$x=-\frac{1}{2},y=\frac{1}{2}$ và $x=2,y=8$.
\end{loi_giai_cau}

\begin{de_cau}{Câu 3}{1,5 điểm}
Một trường tổ chức cho học sinh tham quan. Nếu mỗi xe chở 45 học sinh thì thừa 1 xe và còn 10 chỗ
trống. Nếu mỗi xe chở 40 học sinh thì thiếu 3 xe. Tính số xe và số học sinh.
\end{de_cau}

\begin{loi_giai_cau}
Gọi số xe là $n$, số học sinh là $S$.

- Chở 45/xe: $S=(n-1)\cdot 45+(45-10)=(n-1)\cdot 45+35$.
- Chở 40/xe: cần $n+3$ xe: $S=(n+3)\cdot 40$.

Phương trình: $(n-1)\cdot 45+35=(n+3)\cdot 40$.

$45n-45+35=40n+120\Rightarrow 5n=130\Rightarrow n=26$.

$S=29\cdot 40=1160$ học sinh, $26$ xe.
\end{loi_giai_cau}

\begin{de_cau}{Câu 4}{1,0 điểm}
Cho phương trình $x^2-(2m+1)x+m^2+m-2=0$. Tìm $m$ để phương trình có hai nghiệm $x_1, x_2$
thỏa mãn $x_1-x_2=3$.
\end{de_cau}

\begin{loi_giai_cau}
$\Delta=(2m+1)^2-4(m^2+m-2)=4m^2+4m+1-4m^2-4m+8=9$.

$\Delta=9>0$ với mọi $m$, luôn có 2 nghiệm phân biệt.

$x_1-x_2=\pm\sqrt{\Delta}=\pm 3$.

$|x_1-x_2|=\sqrt{(x_1+x_2)^2-4x_1x_2}=\sqrt{(2m+1)^2-4(m^2+m-2)}=\sqrt{9}=3$.

Luôn thỏa mãn với mọi $m$. Vậy $m\in\mathbb{R}$.
\end{loi_giai_cau}
